\subsection{比较关系：II. 强关系}

定义 3. 给定两个实函数 \(f, g\) ，它们属于 \(\mathcal{H}(\mathfrak{F}, \mathbf{R})\) 且 \(\geqslant 0\) 于 \(\mathfrak{F}\) 的某个集上，我们说 \(f\) 相对于 \(g\) 可忽略，或者说 \(g\) 占优于 \(f\) （沿 \(\mathfrak{F}\) ），并记作 \(f \ll g\) 或 \(g \gg f\) ，如果对每个 \(\varepsilon > 0\) ，存在一个集 \(X \in \mathfrak{F}\) ，使得 \(f(t) \leqslant \varepsilon g(t)\) 对每个 \(t \in X\) 成立。

给定两个赋范空间 \(V_{1}\) 、 \(V_{2}\) 以及两个函数 \(f_{1}\) 、 \(f_{2}\) ，它们分别属于 \(\mathcal{H}(\mathfrak{F}, V_{1})\) 与 \(\mathcal{H}(\mathfrak{F}, V_{2})\) ，我们说 \(f_{1}\) 相对于 \(f_{2}\) 可忽略（沿 \(F\) ），并记作 \(f_{1} \ll f_{2}\) 或 \(f_{2} \gg f_{1}\) ，如果 \(\|f_{1}\| \ll \|f_{2}\|\) 。

对每个标量 \(\lambda \neq 0\) ，关系 \(\mathbf{f}_1 \ll \mathbf{f}_2\lambda\) 等价于 \(\mathbf{f}_1 \ll \mathbf{f}_2\) 。关系 \(\mathbf{f}_1 \ll \mathbf{f}_2\) 蕴含 \(\mathbf{f}_1 \preccurlyeq \mathbf{f}_2\) ，但与之不等价。

注意，关系 \(f_{1} \preccurlyeq f_{2}\) 决不蕴含关系 “ \(f_{1} \ll f_{2}\) 或 \(f_{1} \asymp f_{2}\) ”：有 \(\sin x \preccurlyeq 1\) （当 x 趋于 \(+\infty\) 时），但关系 \(\sin x \ll 1\) 与 \(\sin x \asymp 1\) 中没有一个成立。

例. 1) 关系 \(\mathbf{f} \ll 1\) 表示 \(\mathbf{f}\) 沿 \(\mathfrak{F}\) 趋于 0。

\begin{enumerate}
\def\labelenumi{\arabic{enumi})}
\setcounter{enumi}{1}
\item
  当 \(\alpha\) 与 \(\beta\) 是两个实数且 \(\alpha < \beta\) 时，有 \(x^{\alpha} \ll x^{\beta}\) （当 \(x\) 趋于 \(+\infty\) 时）。类似地，当 \(m\) 与 \(n\) 是两个有理整数且 \(m < n\) 时，有 \(z^{m} \ll z^{n}\) （当复数 \(z\) 趋于 \(\infty\) 时）。
\item
  当 \(x\) 趋于 \(+\infty\) 时，有 \(x^n \ll e^x\) （对每个整数 \(n\) ）(III, p.~105)。
\item
  在 \(R^{2}\) 中，当 \((x, y)\) 趋于 \((0, 0)\) 时，有
\end{enumerate}

\[
x ^ {2} + y ^ {2} \ll | x | + | y |.
\]

下列命题可以由定义 3 立即推出：

命题 4. 若 \(f, g, h\) 是 \(\mathcal{H}(\mathfrak{F}, \mathbf{R})\) 中的三个函数，则关系 \(f \preccurlyeq g\) 与 \(g \ll h\) （相应地， \(f \ll g\) 与 \(g \preccurlyeq h\) ）蕴含 \(f \ll h\) 。

命题 5. 设 \(\mathbf{f}_1, \mathbf{f}_2\) 是 \(\mathcal{H}(\mathfrak{F}, V)\) 中的两个函数， \(g\) 是 \(\mathcal{H}(\mathfrak{F}, \mathbf{R})\) 中的一个函数。则关系 \(\mathbf{f}_1 \ll g\) 与 \(\mathbf{f}_2 \ll g\) 蕴含 \(\mathbf{f}_1 + \mathbf{f}_2 \ll g\) 。

此外，与 V, p.~213 的命题 3 中相同的论证表明：

命题 6. 采用命题 3 的记号，关系 \(\mathbf{f}_1 \ll g_1\) 与 \(\mathbf{f}_2 \preccurlyeq g_2\) （相应地， \(\mathbf{f}_1 \preccurlyeq g_1\) 与 \(\mathbf{f}_2 \ll g_2\) ）蕴含 \([\mathbf{f}_1.\mathbf{f}_2] \ll g_1g_2\) 。

命题 4 表明，关系 \(\mathbf{f} \ll \mathbf{g}\) （ \(\mathcal{H}(\mathfrak{F}, V)\) 中函数之间的）是传递的；但它不是自反的：确切地说，关系 \(\mathbf{f} \ll \mathbf{f}\) 蕴含 \(\mathbf{f}\) 在 \(\mathfrak{F}\) 的某个集上为零（换句话说， \(\mathbf{f}\) 模 \(R_{\infty}\) 等价于 0）；事实上，对一个 \(\varepsilon\) （满足 \(0 < \varepsilon < 1\) ），存在一个 \(X \in \mathfrak{F}\) ，使得 \(\| \mathbf{f}(x) \| \leqslant \varepsilon \| \mathbf{f}(x) \|\) 对每个 \(x \in X\) 成立，而这只有当 \(\mathbf{f}(x) = 0\) 对每个 \(x \in X\) 成立时才可能。由此推出，关系 “ \(\mathbf{f} \ll \mathbf{g}\) 且 \(\mathbf{g} \ll \mathbf{f}\) ” 是传递的和对称的，但不是自反的：因而它不是一个等价关系（它蕴含存在一个集 \(X \in \mathfrak{F}\) ，使得 \(\mathbf{f}(x) = \mathbf{g}(x) = 0\) 对每个 \(x \in X\) 成立）。

命题 7. 若 \(\mathbf{f}\) 与 \(\mathbf{g}\) 是 \(\mathcal{H}(\mathfrak{F}, V)\) 中的两个函数，则关系 \(\mathbf{f} - \mathbf{g} \ll \mathbf{f}\) 等价于 \(\mathbf{f} - \mathbf{g} \ll \mathbf{g}\) 。

事实上， \(\mathbf{f} - \mathbf{g} \ll \mathbf{f}\) 蕴含：对每个 \(\varepsilon > 0\) ，存在一个 \(X \in \mathfrak{F}\) ，使得 \(\| \mathbf{f}(x) - \mathbf{g}(x) \| \leqslant \varepsilon \| \mathbf{f}(x) \|\) 对每个 \(x \in X\) 成立。但这时 \((1 - \varepsilon) \| \mathbf{f}(x) \| \leqslant \| \mathbf{g}(x) \|\) ，从而 \(\mathbf{f} \preccurlyeq \mathbf{g}\) ，由此 (V, p.~215, prop. 4) 得 \(\mathbf{f} - \mathbf{g} \ll \mathbf{g}\) 。

推论. 关系 \(f - g \ll f\) 是 \(\mathcal{H}(\mathfrak{F}, V)\) 上的一个等价关系。

事实上，若 \(f - g \ll f\) 且 \(g - h \ll g\) ，则 \(f - g \ll g\) ，由此 (V, p.~215, prop. 5) 得 \(f - h \ll g\) ，于是，由于 \(g \preccurlyeq f\) ，我们有 \(f - h \ll f\) ，这表明该关系是传递的；它由命题 7 是对称的，且显然是自反的，由此得该推论。

定义 4. 给定 \(\mathcal{H}(\mathfrak{F}, V)\) 中的两个函数 f, g，我们说 f 与 g 等价（沿 F），并记作 \(f \sim g\) ，如果 \(f - g \ll f\) 。

关系 \(\mathbf{f} \sim \mathbf{g}\) 蕴含 \(\mathbf{f} \asymp \mathbf{g}\) ，但与之不等价。

例. 1) 若 \(\mathbf{a}\) 是一个常值函数， \(\neq 0\) 于 \(\mathbf{E}\) 上，则关系 \(\mathbf{f} \sim \mathbf{a}\) 表示 \(\mathbf{f}\) 趋于 \(\mathbf{a}\) （沿 \(\mathfrak{F}\) ）。

\begin{enumerate}
\def\labelenumi{\arabic{enumi})}
\setcounter{enumi}{1}
\item
  设 V 是非离散赋值域 K 上的一个赋范空间，并设 \(\mathbf{f}(x)=\mathbf{a}_{0}x^{n}+\mathbf{a}_{1}x^{n-1}+\cdots+\mathbf{a}_{n}\) 是变量 \(x\in K\) 的一个多项式，其系数在 V 中，且满足 \(a_{0}\neq0\) 。则 \(f(x)\sim a_{0}x^{n}\) （当 \(|x|\) 趋于 \(+\infty\) 时）。
\item
  当实数 \(x\) 趋于 \(+\infty\) 时，有 \(\left(1 + \frac{1}{x}\right)\sin x \sim \sin x\) 。
\item
  当复变量 \(z\) 趋于 0 时，有 \(e^z - 1 \sim z\) 。更一般地，若 \(V\) 是赋值域 \(K\) 上的一个赋范空间， \(\mathbf{f}\) 是定义在 \(x_0 \in K\) 的一个邻域上、取值于 \(V\) 、且具有导数 \(\mathbf{f}'(x_0) \neq 0\) 于点 \(x_0\) 处的函数，则当 \(x\) 趋于 \(x_0\) 时，有 \(\mathbf{f}(x) - \mathbf{f}(x_0) \sim \mathbf{f}'(x_0)(x - x_0)\) (I, p.~3, def. 1)。
\item
  当 \((x, y)\) 趋于 \((0, 0)\) （在 \(R^{2}\) 中）时，有
\end{enumerate}

\[
\sqrt {\sin^ {2} x + \sin^ {2} y} \sim \sqrt {x ^ {2} + y ^ {2}}.
\]

\begin{enumerate}
\def\labelenumi{\arabic{enumi})}
\setcounter{enumi}{5}
\item
  设 \(f(x, y)\) 是两个实变量 x, y 的、以实数为系数的多项式，无常数项。若当 x 保持 \textgreater0 而趋于 0 时，存在一个函数 \(\varphi(x)\) ，它在区间 \([0, a]\) 上连续，且满足 \(\varphi(0) = 0\) 与 \(f(x, \varphi(x)) = 0\) （对 \(0 \leqslant x \leqslant a\) ），则可以证明，存在一个有理数 r 与一个实数 \(\lambda \neq 0\) ，使得 \(\varphi(x) \sim \lambda x^{r}\) (V, p.~259, exerc. 3)。
\item
  对每个 \(x > 0\) ，用 \(\pi(x)\) 表示 \(\leqslant x\) 的素数的个数；已经证明， \(\pi(x) \sim x / \log x\) （当 \(x\) 趋于 \(+\infty\) 时）。
\end{enumerate}

\pandocbounded{\includegraphics[keepaspectratio]{images/part002_abf04123.jpg}}

注. 注意，关系 \(\mathbf{f} \sim \mathbf{g}\) 决不蕴含差 \(\mathbf{f} - \mathbf{g}\) 沿 \(\mathfrak{F}\) 趋于 0；这个差甚至可以无界，正如例子 \(x^{2} + x \sim x^{2}\) （当 \(x\) 趋于 \(+\infty\) 时）所表明的。

命题 8. 设 K 是一个赋值域， \(f_{1}, f_{2}, g_{1}, g_{2}\) 是 \(\mathcal{H}(\mathfrak{F}, \mathrm{K})\) 中的四个函数；则关系 \(f_{1} \sim g_{1}\) 与 \(f_{2} \sim g_{2}\) 蕴含 \(f_{1}f_{2} \sim g_{1}g_{2}\) 。

事实上，我们有 \(f_{1}f_{2}-g_{1}g_{2}=f_{1}(f_{2}-g_{2})+(f_{1}-g_{1})g_{2}\) ；由于 \(f_{1}\preccurlyeq g_{1}\) 、 \(f_{1}-g_{1}\ll g_{1}\) 与 \(f_{2}-g_{2}\ll g_{2}\) ，我们有 \(f_{1}f_{2}-g_{1}g_{2}\ll g_{1}g_{2}\) (V, p.~215, prop. 5 and 6)。

相反，我们在 V, p.~214 中已给出一个例子，其中有 \(f_{1} = g_{1}\) 、 \(f_{2} \sim g_{2}\) ，然而关系 \(f_{1} + f_{2} \asymp g_{1} + g_{2}\) 并不成立（故更谈不上

\[
f _ {1} + f _ {2} \sim g _ {1} + g _ {2}).
\]

比较关系 \(f \ll g\) 、 \(f \sim g\) 称为强关系。 \(\mathcal{H}(\mathfrak{F}, V)\) 中的两个函数 f, g 称为可比较的（当想要避免任何可能的混淆时，称为强可比较的），如果它们满足下列三个关系之一： \(f \ll g\) 、 \(f \gg g\) ，或 “存在一个 \(\lambda \neq 0\) 使得 \(f \sim g\lambda\) ”。

注. 1) 两个函数可以是弱可比较的但不是强可比较的，例如函数 1 与 \(\sin x\) （当 \(x\) 趋于 \(+\infty\) 时）。

\begin{enumerate}
\def\labelenumi{\arabic{enumi})}
\setcounter{enumi}{1}
\tightlist
\item
  在比较关系 \(f_{1} \preccurlyeq f_{2}\) 与 \(f_{1} \ll f_{2}\) 的定义中，空间 \(V_{1}, V_{2}\) （ \(f_{1}\) 与 \(f_{2}\) 分别在其中取值）上的范数只是表面上介入；实际上介入的只是 \(V_{1}\) 与 \(V_{2}\) 上的拓扑，因为关系 \(f_{1} \preccurlyeq f_{2}\) 与 \(f_{1} \ll f_{2}\) 当把 \(V_{1}\) 或 \(V_{2}\) 上的范数换成一个等价范数时，就被等价的关系所代替 (Gen.~Top., IX, p.~170, def. 7)。
\end{enumerate}

